% -*- TeX:RNW:UK -*-
\documentclass[10pt]{beamer}
\usetheme{metropolis}
%\useinnertheme{rectangles}
\setbeamercovered{%
still covered={\opaqueness<1->{15}},
again covered={\opaqueness<1->{40}}}

\hypersetup{colorlinks,linkcolor=black,urlcolor=brown,citecolor=brown}

\usepackage{amsmath,amssymb,amsthm}
\usepackage{unicode-math}

% We set the Lucida OTF fonts as default
\usepackage{fontspec}
\setmainfont{Lucida Bright OT}
\setsansfont{Lucida Sans OT}
\setmonofont{Lucida Console DK}[Scale=MatchLowercase]

\newfontfamily\typiconsfont{Typicons}[Scale=0.9]
\newcommand\typicons[1]{{\typiconsfont\symbol{#1}}}
\newcommand\Attention{\colorbox{yellow}{\textcolor{black}{\typicons{"E136}}}\xspace}
\newcommand\Advanced{\colorbox{black}{\textcolor{white}{\typicons{"E04E}}}\xspace}

\newfontfamily\webglyphsfont{WebHostingHub-Glyphs}[Scale=0.7]
\newcommand\webglyphs[1]{{\webglyphsfont\symbol{#1}}}
\newcommand\Discussion{\colorbox{white}{\textcolor{black}{\webglyphs{"F134}}}\xspace}
\newcommand\DExamples{\colorbox{black}{\textcolor{white}{\webglyphs{"F134} examples?}}}
\newcommand\Reading{\colorbox{black}{\textcolor{white}{\webglyphs{"F0C1}}}\xspace}

\newfontfamily\meteoconsfont{Meteocons}
\newcommand\meteocons[1]{{\meteoconsfont\symbol{#1}}}
\newcommand\meteosun{\meteocons{"0042}}
\newcommand\meteosuncloud{\meteocons{"0048}}
\newcommand\meteorain{\meteocons{"0052}}
\newcommand\meteosolidsun{\meteocons{"0031}}

\newfontfamily\KRfarmfont{KR Back On The Farm}
\newcommand\KRfarm[1]{{\KRfarmfont\symbol{#1}}}
\newcommand\farmplant{\KRfarm{"0049}}

\usepackage{polyglossia}
\setdefaultlanguage[variant = british, ordinalmonthday = false]{english}

\usepackage[style=authoryear-comp,firstinits,sortcites,maxcitenames=2,%
    mincitenames=1,maxbibnames=10,minbibnames=10,uniquename=mininit,%
    uniquelist=minyear,sortfirstinits=true]{biblatex}
\addbibresource{../references/ecophys.bib}
%\renewcommand{\bibfont}{\small}

\usepackage{abbrev}
\DeclareGraphicsExtensions{.png,.jpg,.jpeg,.pdf}

\begin{document}

\title{PBIO-141\\Sensory and Physiological Ecology of  Plants}
\subtitle{10: Carbon Assimilation (A)}
\author{Pedro J. Aphalo}
\date{January-February 2018}
\institute[Univ.\ of Helsinki]{M.Sc.\ in Plant Biology, University of Helsinki\\[2ex] \url{http://blogs.helsinki.fi/aphalo/}}


  \begin{frame}
    \maketitle
  \end{frame}

  \begin{frame}[c]
    \begin{center}
      \begin{small}
        \copyright 2006-2018 by Pedro J. Aphalo\\
       University of Helsinki, Finland.\\
        \textcolor{blue}{\url{http://blogs.helsinki.fi/aphalo/}}\\[2ex]
      \end{small}

      \begin{footnotesize}
        Sensory and Physiological Ecology of  Plants: `10: Carbon Assimilation (A)' by Pedro J. Aphalo is licensed under a Creative Commons Attribution-ShareAlike 4.0 International License.

        Illustrations and text quoted from copyrighted sources is excluded from this license and their use should respect the original licenses.\\[2ex]
      \end{footnotesize}

      \includegraphics[width=6em]{../figures/copyright/by-sa}
    \end{center}
  \end{frame}

  \begin{frame}
    \frametitle{Outline}
    \tableofcontents
  \end{frame}

\section{Carbon assimilation}

\begin{frame}{Carbon economy}
    \begin{itemize}
        \item Our objective is to describe and understand acclimation, adaptation
        and responses of carbon acquisition and loss by whole
        plants.
        \item For understanding we need mechanistic explanations.%
        \item So, we need to study plants at different scales.%
        \item From biochemistry upwards.
    \end{itemize}
\end{frame}

\begin{frame}{Scale}
    \centering
    \includegraphics[height=0.8\textheight]{figures/OsmCho1988-1}\\
    {\small Carbon assimilation and growth at different scales
    \autocite[from][]{OsmCho1988}.}
\end{frame}

\begin{frame}{Forest}
    \centering
    \includegraphics[width=0.7\textwidth]{figures/Koli}\\
    {\small Koli national park in summer. Growth and changes in the forest
    can be observed over a period of years. The relaxation time after a
    disturbance is long.}
\end{frame}

\begin{frame}{Single plant}
    \centering
    \includegraphics[height=0.70\textheight]{figures/Linnea-1}\\
    {\small Twinflower (vanamo) (\species{Linnea borealis}). The relaxation time after a disturbance is
    shorter for a single plant than for the forest.}
\end{frame}

\begin{frame}{Leaf cross-section}
    \centering
    \includegraphics[width=0.7\textwidth]{figures/IvyCryofractureSEM-1}\\
    {\small Cross-section of an ivy (muratti) leaf (\species{Hedera helix}), cryofracture SEM \autocite[from][]{Aphalo1991}.}
\end{frame}

\begin{frame}{Chloroplast ultrastructure}
    \centering
    \includegraphics[width=0.8\textwidth]{figures/TaizXXChloroplastTEM-1}\\
    {\small Chloroplast cross-section, TEM
    \autocite[from][]{TaiZei2006}.}
\end{frame}

\begin{frame}{Carbon assimilation}
    {Discussion: How does C assimilation take place? (Whole path from atmosphere to carbohydrate.)}

    \begin{itemize}
        \item Approximately 40\% of the plant dry mass consists of carbon.

        \item Life on earth depends on photosynthesis.

        \item \COtwo uptake occurs through leaf pores, the stomata, which are able
        to rapidly change their aperture.

        \item Once inside the leaf, \COtwo diffuses from the intercellular air
        spaces to the site of carboxylation.

        \item The primary processes of photosynthesis occur in the chloroplast.
    \end{itemize}
\end{frame}

\begin{frame}{Gas-exchange in a leaf}
    \centering
    \includegraphics[height=0.7\textheight]{figures/Heldt8-1LeafStoma}\\
    {\small Gas-exchange through stomata in a leaf \autocite[from][]{Heldt1997}.}
\end{frame}

\begin{frame}{The leaf}
    \begin{itemize}
        \item In plant leaves respiration takes place at all times,
        and photosynthesis when there is light.
        \item The ``real'' rate of carbon fixation is called gross
        photosynthetic rate.
        \item The gross photosynthetic rate less respiration
        gives the net photosynthetic rate.
        \item When we measure the exchange of \COtwo\ between a leaf
        and the atmosphere we measure the net photosynthetic rate.
    \end{itemize}
    {$$ A_N = A_G - R $$
    where $A$ is the usual symbol for photosynthesis (=assimilation)
    and $R$ the usual symbol for respiration.}
\end{frame}

\begin{frame}{Measuring leaf gas-exchange}
    \centering
    \includegraphics[height=0.7\textheight]{figures/licor}\\
    {\small Photo source: \url{http://www.geos.ed.ac.uk/homes/s0452637/ecophys.html}}
\end{frame}


\section{Photosynthesis at other scales}

\begin{frame}{Big leaf models}
    \begin{itemize}
        \item Sometimes a canopy is treated as a \emph{big leaf}.
        \item $A$ and $R$ are measured and expressed per unit
        \emph{ground} area.
        \item Scaling from leaf measurements to `big leaf' values is
        a complicated and error prone process.
        \item The gas exchange of canopies can also be measured directly.
    \end{itemize}
\end{frame}

\begin{frame}{Eddy covariance}
    \centering
    \includegraphics[width=0.8\textwidth]{figures/EddyCovariance_diagram_1}\\
    {\small Representation of eddies passing a measuremnet tower. Diagram created by George Burba.}
    {\tiny \url{http://en.wikipedia.org/wiki/File:EddyCovariance_diagram_1.jpg}}
\end{frame}

\begin{frame}{Eddy covariance}
    \centering
    \includegraphics[height=0.7\textheight]{figures/BobCook5_8Apr00}\\
    {\small Tower used for eddy covariance measurements} {\tiny \url{http://www.nrel.colostate.edu/projects/kruger/Images/BobCook5_8Apr00.jpg}}
\end{frame}

\begin{frame}{Chloroplasts}
    \begin{itemize}
        \item Thought to have originated very long ago during
        evolution from an endosymbiont.
        \item Chloroplasts have their own genetic information, and
        multiply by division.
        \item Only a fraction of the proteins and enzymes in the
        chloroplast are coded by chloroplast genes, the rest are
        encoded by nuclear genes.
        \item Plant physiologists sometimes study and do experiments
        with isolated chloroplasts (chloroplasts separated from the
        cells).
    \end{itemize}
\end{frame}

\begin{frame}{Biochemistry}
    Not all plants have the same type of photosynthetic
    metabolism.
    \begin{description}
        \item[C3] The first type described. The first product is
        a molecule with three carbon atoms.
        \item[C4] The first product is a molecule with four carbon
        atoms. Its pathway includes a \COtwo\ concentrating stage.
        There are several subtypes. Appeared
        more than once during evolution of plants.
        \item[CAM] It includes a \COtwo\ ``storage'' stage.
        Accumulation occurs at night, reduction and synthesis of
        carbohydrates during the day.
    \end{description}
\end{frame}

\section{Biochemistry: C3}

\begin{frame}{C3 photosynthesis}
    \begin{itemize}
        \item The C3 metabolism is the simplest one, and it is also a
        component of C4 metabolism and CAM.
        \item All the steps can happen inside a single chloroplast.
        \item It is also the most widespread type of carbon assimilation
        metabolism in the plant kingdom.
    \end{itemize}
\end{frame}

\begin{frame}{C3 photosynthesis}
    \begin{itemize}
        \item In photosynthesis energy from light is used to
        assimilate C from \COtwo\ in the air (or water).
        \item C is not only fixed by carboxilation but also reduced.
        \item Triose phosphate is produced and further used to
        synthesize carbohydrates (sugars, starch).
        \item The name \emph{photosynthate} is used as a generic
        name for these products of photosynthesis.
        \item Starch can accumulate in the chloroplasts themselves,
        or carbohydrates may be transported away from the leaves and
        used elsewhere in the plant as a source of building blocks
        and energy for growth or maintenance.
    \end{itemize}
\end{frame}

\begin{frame}{C3 photosynthesis}
    \centering
    \includegraphics[height=0.65\textheight]{figures/Heldt6-1PhotosynBasic}\\
    {\small Basic scheme \autocite[from][]{Heldt1997}.}
\end{frame}

\begin{frame}{Light and carbon reactions}
    \begin{itemize}
        \item C3 photosynthesis can be divided into two stages:
        \emph{light reactions} and \emph{carbon reactions}.
        \item[$\circ$] Earlier the name \emph{dark reactions} was
        used for what now are called carbon reactions.
        \item Light reactions are related to capturing light energy
        into energy rich ATP and reducing power as NADPH.
        \item Carbon reactions are related to fixing and reducing C
        from \COtwo. They use as substrate the products of the light
        reactions. They can proceed in darkness if ATP and NADPH are
        supplied.
    \end{itemize}
\end{frame}

\begin{frame}{Light and carbon reactions}
    \centering
    \includegraphics[width=\textwidth]{figures/Taiz8-1ChloroplastLightAndCarbonReactions}\\
    {\small Location of light and carbon reactions in the chloroplast \autocite[from][]{TaiZei2006}.}
\end{frame}

\section{Chloroplast}

\begin{frame}{Chloroplast: location of photosynthesis}
    \centering
    \includegraphics[width=1.1\textwidth]{figures/Taiz7-16Chloroplast}\\
    {\small Diagram of a chloroplast \autocite[from][]{TaiZei2006}.}
\end{frame}

\section{C3: Light reactions}

\begin{frame}{Photosystems}
    \begin{itemize}
        \item Two types of photosystems are involved in the light
        reactions: PS I and PS II.
        \item They play different roles, and have at their centre
        chlorophylls that absorb light of slightly different
        wavelengths.
        \item Each PS is composed of many pigment molecules and
        several proteins.
        \item There are different regions: the antenna complexes
        and the reaction centres.
    \end{itemize}
\end{frame}

\begin{frame}{Light reactions}
    \begin{itemize}
        \item Chlorophyll and accessory pigments (mainly
        carotenoids) absorb light.
        \item When one of these molecules absorbs a photon, it
        becomes excited.
        \item This excitation can be the start of a
        photochemical reaction, or it can be transferred to another
        molecule. (It can also be dissipated as heat or emitted as
        light of a longer wavelength.)
    \end{itemize}
\end{frame}

\begin{frame}{Reaction centres}
    \begin{itemize}
        \item Chlorophyll molecules are not organized at random in
        the chloroplast.
        \item There are reaction centres and antenna complexes.
        \item The excitation of chlorophyll molecules in the
        antenna complexes does not drive directly the light
        reactions.
        \item The excitation is funnelled to special chlorophylls at
        the reaction centres.
        \item The photochemical reactions happen in these reaction
        centres.
    \end{itemize}
\end{frame}

\begin{frame}{Light reactions}
    \centering
    \includegraphics[height=0.85\textheight]{figures/Heldt3-17PhotosyntChloroplast}\\
    {\small Simplified scheme of light reactions \autocite[from][]{Heldt1997}.}
\end{frame}

\begin{frame}{Physical organization}
    \begin{itemize}
        \item PS I and PS II are not located side by side in the
        membrane, they are located in different parts of the
        lamellae.
        \item PS I in stroma lamellae.
        \item PS II in grana lamellae.
        \item The electron carriers move to transfer the electrons.
    \end{itemize}
\end{frame}

\begin{frame}{Pigments}
    \centering
    \includegraphics[height=0.85\textheight]{figures/Heldt2-3ChlCarot}\\
    {\small Absorption spectra and spectral (energy) irradiance\\ \autocite[from][]{Heldt1997}.}
\end{frame}

\begin{frame}{Funneling}
    \centering
    \includegraphics[width=0.9\textwidth]{figures/Taiz7-19Funneling}\\
    {\small Funneling of energy into the reaction centers \autocite[from][]{TaiZei2006}.}
\end{frame}

%\begin{frame}{Electron transport chain}
%    \begin{itemize}
%        \item Electrons from the water molecule split at PS II are
%        transported through the electron transport chain to the
%        final acceptor NADP$^+$ which is reduced to NADPH.
%        \item The chain is located in the thylakoid membrane of the
%        chloroplasts but it is not made up of fixed directly linked
%        components.
%        \item Different parts of the chain are located in different
%        regions of the membrane and some of the intermediate acceptors like
%        plastocyanine are soluble and move between the regions transferring
%        the electrons.
%    \end{itemize}
%\end{frame}

%\begin{frame}{Z scheme}
%    \vspace{-0.04\textheight}
%    \includegraphics[width=\textwidth]{figures/Taiz7.14Zscheme}\\
%%    \vspace{-0.02\textheight}
%    {\small Z scheme \autocite[from][]{TaiZei2006}.}
%\end{frame}

\begin{frame}{Discussion (5 min)}
  Can you suggest ways in which light reaction `machinery' could acclimate to shade?
\end{frame}

\section{C3: Carbon reactions}

\begin{frame}{Carbon reactions}
\begin{itemize}
        \item In the carbon reactions ATP (adenosine triphosphate) and NADPH
        (nicotinamide adenine dinucleotide phosphate) from the
        light reactions are used.
        \item The product is a triose phosphate.
        \item A \emph{tri}ose is a three carbon monosaccharide (sugar).
    \end{itemize}
\end{frame}

\begin{frame}{Carbon reactions}
    \centering
    \includegraphics[height=0.85\textheight]{figures/Heldt6-2PhotosynOverall}\\
    {\small Overall ``dark'' or carbon reaction \autocite[from][]{Heldt1997}.}
\end{frame}

\begin{frame}{Carbon reactions}
    \begin{itemize}
        \item Carbon reactions take place as a cycle.
        \item This cycle is called Calvin cycle.
        \item The key enzyme in the cycle is Rubisco.
        \item Rubisco is a short name for ribulose 1,5-bisphosphate
        carboxylase oxygenase.
        \item As the name implies it can both carboxylate and oxygenate
        its substrate. Obviously, only carboxylation leads to carbon
        assimilation.
        \item The double function of Rubisco means that there is
        competition between two reactions with the partial rates
        depending on the concentrations of \COtwo\ and O$_2$.
        \item The affinity of Rubisco for \COtwo\ relative to O$_2$ decreases
        with increasing temperature.
    \end{itemize}
\end{frame}

\begin{frame}{Melvin Calvin}
    \begin{itemize}
        \item The name \emph{Calvin cycle} is used for the
        reductive pentose phosphate pathway.
        \item Melvin Calvin described it and received the
        Nobel prize in 1961. This discovery was done using $^{14}$C to trace
        the fate of the C in \COtwo assimilated. There were other
        people in his team contributing to the research.
    \end{itemize}
\end{frame}

%\begin{frame}{Calvin cycle}
%    \vspace{-0.04\textheight}
%    \includegraphics[height=0.65\textheight]{figures/Heldt6.3CalvinOverview}\\
%%    \vspace{-0.02\textheight}
%    {\small Overview of Calvin cycle \autocite[from][]{Heldt1997}.}
%\end{frame}
%
\begin{frame}{Calvin cycle}
    \centering
    \includegraphics[width=0.85\textwidth]{figures/Taiz8-2CalvinStages}\\
    {\small The three stages of the Calvin cycle: carboxylation, reduction and regeneration \autocite[from][]{TaiZei2006}.}
\end{frame}
%
%\begin{frame}{Calvin cycle: carboxylation}
%    \begin{itemize}
%        \item Ribulose is a pentose.
%        \item A pentose is a five carbon monosaccharide.
%        \item Ribulose 1,5-bisphosphate is phosphorylated at
%        positions 1 and 5.
%        \item After carboxylation it yields two molecules of
%        3-P-glycerate.
%        \item 3-P-glycerate is a three-carbon molecule.
%        \item So we have a total of six carbon atoms, five from the
%        ribulose and one from \COtwo.
%%        \item<7> To get one triose phosphate as product we have to
%%        start with three ribulose 1,5-bisphosphate molecules and
%%        three \COtwo molecules.
%%        \item<8> Water is an additional substrate in the carboxylation reactions.
%    \end{itemize}
%\end{frame}
%
%\begin{frame}{Calvin cycle: carboxylation}
%    \begin{itemize}
%        \item To get one triose phosphate as product we have to
%        start with three ribulose 1,5-bisphosphate molecules and
%        three \COtwo molecules.
%        \item Water is an additional substrate in the carboxylation reactions.
%    \end{itemize}
%\end{frame}

%\begin{frame}{Calvin cycle}
%%    \vspace{-0.06\textheight}
%    \twocolumn{%
%    \includegraphics[height=\textheight]{figures/Taiz8.3Calvin}}{%
%    \raggedright
%    \vspace{0.9\textheight}
%    {\small Detail of the Calvin cycle \autocite[from][]{TaiZei2006}.}}
%\end{frame}
%
%\begin{frame}{Calvin cycle: reduction}
%    \begin{itemize}
%        \item The 3-P-glycerate is reduced to
%        glyceraldehyde-3-phosphate.
%        \item The reaction uses NADPH as a source of
%        reducing power (``a source of H'') and ATP as a source of P,
%        removing one oxygen atom from the 3-P-glycerate as PO$_4H^{2-}$
%        (phosphate) and yielding glyceraldehyde-3-P.
%    \end{itemize}
%\end{frame}
%
%\begin{frame}{Calvin cycle: regeneration}
%    \begin{itemize}
%        \item This is the regeneration of three ribulose 1,5-bisphosphate
%        molecules from five triose phosphate molecules (the sixth
%        triose phosphate molecule is the product of going through the cycle once).
%        \item Three ribulose phosphate molecules are first produced
%        and then phosphorylated to bisphosphate using ATP in the
%        process.
%        \item[$\circ$] (The sixth triose phosphate molecule is
%        exported to the cytosol.)
%    \end{itemize}
%\end{frame}
%
%\begin{frame}{Oxygenation}
%    \begin{itemize}
%        \item Rubisco catalyzes two reactions: the carboxylation of
%        ribulose 1,5-bisphosphate or its oxygenation.
%        \item These reactions compete, and their relative reaction
%        rates will depend at a given temperature on the concentrations of \COtwo\ and
%        oxygen.
%        \item The relative affinity of Rubisco for \COtwo and O$_2$
%        depends on temperature, with high temperatures favouring
%        oxygenation.
%        \item Oxygenation gives rise to photorespiration, which
%        releases \COtwo.
%    \end{itemize}
%\end{frame}
%
%\begin{frame}{Photorespiration}
%    \begin{itemize}
%        \item The products of the oxygenation of ribulose
%        1,5-bisphosphate are 1 molecule of 2-phosphoglycolate (2
%        carbons) and 3-phosphoglycerate.
%        \item 2-phosphoglycolate is recycled but at a rather high
%        cost in ATP and NADPH.
%        \item The oxygenase side reaction can cost to the plant more
%        than one third of the captured photons.
%    \end{itemize}
%\end{frame}
%
%\begin{frame}{Photorespiration}
%    \begin{itemize}
%        \item The role of photorespiration is not well understood,
%        but it reduces the efficiency of photosynthesis, especially
%        at high temperature.
%        \item On the other hand it could work as a safety valve
%        when under high light and low \COtwo (closed stomata) there
%        is an excess of ATP and NADPH produced in the light reactions.
%        \item C4 metabolism ensures that carboxylation
%        predominates, even at high temperature, but it has an
%        energetic cost.
%    \end{itemize}
%\end{frame}

%\begin{frame}{Photosystems}
%    \begin{itemize}
%        \item Two types of photosystems are involved in the light
%        reactions: PS I and PS II.
%        \item They play different roles, and have at their centre
%        chlorophylls that absorb light of slightly different
%        wavelengths.
%        \item Each PS is composed of many pigment molecules and
%        several proteins.
%        \item There are different regions: the antenna complexes
%        and the reaction centres.
%    \end{itemize}
%\end{frame}
%
%\begin{frame}{Light reactions}
%    \begin{itemize}
%        \item Chlorophyll and accessory pigments (mainly
%        carotenoids) absorb light.
%        \item When one of these molecules absorbs a photon, it
%        becomes excited.
%        \item This excitation can be the start of a
%        photochemical reaction, or it can be transferred to another
%        molecule. (It can also be dissipated as heat or emitted as
%        light of a longer wavelength.)
%    \end{itemize}
%\end{frame}
%
%\begin{frame}{Reaction centres}
%    \begin{itemize}
%        \item Chlorophyll molecules are not organized at random in
%        the chloroplast.
%        \item There are reaction centres and antenna complexes.
%        \item The excitation of chlorophyll molecules in the
%        antenna complexes does not drive directly the light
%        reactions.
%        \item The excitation is funnelled to special chlorophylls at
%        the reaction centres.
%        \item The photochemical reactions happen in these reaction
%        centres.
%    \end{itemize}
%\end{frame}
%
%\begin{frame}{Light reactions}
%    \vspace{-0.07\textheight}
%    \includegraphics[height=0.8\textheight]{figures/Heldt3.17PhotosyntChloroplast}\\
%%    \vspace{-0.02\textheight}
%    {\small Simplified scheme of light reactions \autocite[from][]{Heldt1997}.}
%\end{frame}
%
%\begin{frame}{Physical organization}
%    \begin{itemize}
%        \item PS I and PS II are not located side by side in the
%        membrane, they are located in different parts of the
%        lamellae.
%        \item PS I in stroma lamellae.
%        \item PS II in grana lamellae.
%        \item The electron carriers move to transfer the electrons.
%    \end{itemize}
%\end{frame}
%
%\begin{frame}{Chloroplast}
%    \vspace{-0.04\textheight}
%    \includegraphics[width=1.1\textwidth]{figures/Taiz7.16Chloroplast}\\
%%    \vspace{-0.02\textheight}
%    {\small Diagram of a chloroplast \autocite[from][]{TaiZei2006}.}
%\end{frame}
%
%\begin{frame}{Pigments}
%    \vspace{-0.04\textheight}
%    \includegraphics[height=0.65\textheight]{figures/Heldt2.3ChlCarot}\\
%%    \vspace{-0.02\textheight}
%    {\small Absorption spectra \autocite[from][]{Heldt1997}.}
%\end{frame}
%
%\begin{frame}{Funneling}
%    \vspace{-0.04\textheight}
%    \includegraphics[width=\textwidth]{figures/Taiz7.19Funneling}\\
%%    \vspace{-0.02\textheight}
%    {\small Funneling of energy into the reaction centers \autocite[from][]{TaiZei2006}.}
%\end{frame}
%
%\begin{frame}{Electron transport chain}
%    \begin{itemize}
%        \item Electrons from the water molecule split at PS II are
%        transported through the electron transport chain to the
%        final acceptor NADP$^+$ which is reduced to NADPH.
%        \item The chain is located in the thylakoid membrane of the
%        chloroplasts but it is not made up of fixed directly linked
%        components.
%        \item Different parts of the chain are located in different
%        regions of the membrane and some of the intermediate acceptors like
%        plastocyanine are soluble and move between the regions transferring
%        the electrons.
%    \end{itemize}
%\end{frame}
%
%\begin{frame}{Z scheme}
%    \vspace{-0.04\textheight}
%    \includegraphics[width=\textwidth]{figures/Taiz7.14Zscheme}\\
%%    \vspace{-0.02\textheight}
%    {\small Z scheme \autocite[from][]{TaiZei2006}.}
%\end{frame}
%
%\begin{frame}{Next}
%    \begin{itemize}
%        \item C3, carbon reactions
%        \item C4
%        \item CAM
%        \item[]
%        \item Whole-leaf response curves
%        \item Acclimation and adaptation of photosynthesis
%    \end{itemize}
%\end{frame}
%
%\begin{frame}{Carbon reactions}
%    \begin{itemize}
%        \item In the carbon reactions ATP and NADPH from the
%        light reactions is used.
%        \item The product is a triose phosphate.
%        \item A \emph{tri}ose is a three carbon sugar.
%    \end{itemize}
%\end{frame}
%
%\begin{frame}{Carbon reactions}
%    \vspace{-0.04\textheight}
%    \includegraphics[height=0.65\textheight]{figures/Heldt6.2PhotosynOverall}\\
%%    \vspace{-0.02\textheight}
%    {\small Overall ``dark'' or carbon reaction \autocite[from][]{Heldt1997}.}
%\end{frame}

  \section{References}
  
  \begin{frame}[t,allowframebreaks]{References}
    \nocite{Heldt2004}
    \printbibliography
  \end{frame}

\end{document}

